% Part: first-order-logic
% Chapter: model-theory
% Section: substructures

\documentclass[../../../include/open-logic-section]{subfiles}

\begin{document}

\olfileid{mod}{bas}{sub}
\olsection{فرعي \printtoken{p}{structure}}

د يوه !!a{structure}~$\Struct{M}$ !!{domain} ښايي د بل
جوړښت~$\Struct{M'}$ فرعي سټ وي۔ خو څرګنده ده چې $\Struct{M}$ بايد
يوازې هغه وخت د~$\Struct{M'}$ يوه ``برخه'' وګڼو چې نۀ يوازې
$\Domain{M} \subseteq \Domain{M'}$ وي، بلکې $\Struct{M}$ او
$\Struct{M'}$ لږ تر لږه پر ګډې برخې~$\Domain{M}$ د ژبې د نښو په
تفسيرولو کښې هم ``سره سلا'' وي۔

\begin{defn}
\ollabel{defn:substructure}
د يوې ژبې~$\Lang L$ !!{structure}s $\Struct M$ او $\Struct M'$ راکړل شي۔
وايو $\Struct M$ د~$\Struct M'$ \emph{فرعي !!{structure}}، او
$\Struct M'$ د~$\Struct M$ \emph{غځونه} ده، چې
$\Struct M \substruct \Struct M'$ ئې ليکو، هله او يوازې هله چې
\begin{enumerate}
\item $\Domain{M} \subseteq \Domain{M'}$ وي؛
\item د هر داسې ثابت $c \in \Lang L$ دپاره، $\Assign{c}{M} =
    \Assign{c}{M'}$ وي؛
\item د هر داسې $n$-ځاييز !!{function}~$f \in \Lang L$ دپاره، د
  ټولو $a_1$، \dots،~$a_n \in \Domain{M}$ دپاره
  $\Assign{f}{M}(a_1, \dots, a_n) = \Assign{f}{M'}(a_1, \dots, a_n)$
  وي؛
\item د هر داسې $n$-ځاييز !!{predicate}~$R \in \Lang L$ دپاره، د
  ټولو $a_1$، \dots،~$a_n \in \Domain{M}$ دپاره
  $\langle a_1, \dots, a_n\rangle \in \Assign{R}{M}$ هله او يوازې
  هله چې $\langle a_1, \dots, a_n\rangle \in \Assign{R}{M'}$ وي۔
\end{enumerate}
\end{defn}

\begin{rem}
\ollabel{rem:substructure}
که ژبه ثابت يا !!{function}s ونۀ لري، نو هر
$N \subseteq \Domain{M}$ د~$\Struct M$ يو فرعي
!!{structure}~$\Struct{N}$ ټاکي چې !!{domain} ئې
$\Domain{N} = N$ دے؛ د دې دپاره $\Assign{R}{N} =
\Assign{R}{M} \cap N^n$ ږدو۔
\end{rem}

% prove something about this? Examples?

\end{document}
